% Generated from locale/tr/content/first-order-logic/beyond/intuitionistic-logic.tex; do not hand-edit.


\olfileid[tr]{fol}{byd}{il}

\olsection{Sezgici Mantık}

İkinci dereceden ve yüksek dereceden mantığın tersine, sezgici birinci
dereceden mantık klasik sürümün bir kısıtlamasıdır ve daha ``yapıcı'' bir akıl
yürütme türünü modellemeyi amaçlar. Aşağıdaki örnekler, bunun altında yatan
bazı güdüleri açıklamaya yarayabilir.

Birinin bir gün yanınıza gelip şu özelliğe sahip bir $x$ doğal sayısı
belirlediğini açıkladığını varsayın: $x$ asal ise Riemann hipotezi doğru, $x$
bileşik ise Riemann hipotezi yanlıştır. Müthiş haber!{} Riemann hipotezinin
doğru olup olmadığı matematiğin büyük açık sorularından biridir; oysa bu kişi
sorunu bir hesaplama problemine, yani belirli bir sayının asal olup olmadığını
belirlemeye indirgemiş görünmektedir.

$x$'in sihirli değeri nedir? Şöyle betimliyor: $x$, Riemann hipotezi doğruysa
$7$'ye, değilse $9$'a eşit olan doğal sayıdır.

Öfkelenip paranızı geri istersiniz. Klasik bakımdan yukarıdaki betimleme
gerçekten de $x$'in tek bir değerini belirler; ama sizin asıl istediğiniz,
\emph{açıkça} verilen bir $x$ değeridir.

Başka, belki daha az yapay bir örnek olarak şu soruyu ele alın. İrrasyonel bir
sayının rasyonel bir kuvvetini alıp rasyonel bir sonuç elde etmenin mümkün
olduğunu biliyoruz. Örneğin $\sqrt{2}^2 = 2$. Daha az açık olan, irrasyonel bir
sayının \emph{irrasyonel} bir kuvvetini alıp rasyonel bir sonuç elde etmenin
mümkün olup olmadığıdır. Aşağıdaki teorem buna olumlu yanıt verir:

\begin{thm}
İrrasyonel $a$ ve $b$ sayıları vardır ve $a^b$ rasyoneldir.
\end{thm}

\begin{proof}
$\sqrt{2}^{\sqrt{2}}$ sayısını ele alın. Bu sayı rasyonelse işimiz bitmiştir:
$a = b = \sqrt{2}$ alabiliriz. Değilse irrasyoneldir. O zaman
\[
(\sqrt{2}^{\sqrt{2}})^{\sqrt{2}} = \sqrt{2}^{\sqrt{2} \cdot
  \sqrt{2}} = \sqrt{2}^2 = 2,
\]
olur ve bu kesinlikle rasyoneldir. Dolayısıyla bu durumda $a$'yı
$\sqrt{2}^{\sqrt{2}}$, $b$'yi de~$\sqrt 2$ alın.
\end{proof}

Bu, geçerli bir kanıt sayılır mı? Çoğu matematikçi öyle olduğunu düşünür.
Ancak burada da biraz doyurucu olmayan bir şey vardır: belirli bir özelliğe
sahip bir gerçek sayı çiftinin varlığını, bunun \emph{hangi} sayı çifti
olduğunu söyleyemeden kanıtladık. Aynı sonuç, $a$, $b$ çifti kanıtta
\emph{verilecek} biçimde de kanıtlanabilir: $a = \sqrt{3}$ ve $b = \log_3 4$
alın. O zaman
\[
a^b = \sqrt{3}^{\log_3 4} = 3^{1/2 \cdot \log_3 4} = (3^{\log_3
  4})^{1/2} = 4^{1/2}= 2,
\]
elde edilir; çünkü $3^{\log_3 x} = x$.

Sezgici mantık, ilk kanıttaki gibi hamlelere izin verilmeyen bir akıl yürütme
türünü modellemek üzere tasarlanmıştır. Bir $x$ bulup onun $!A(x)$ koşulunu
sağladığını göstererek bir varlık savını kanıtlamak, ikinci kanıttaki gibi
belirli bir~$x$ ile onun $!A$
koşulunu sağladığının bir kanıtını vermeniz gerektiği anlamına gelir. $!A$ ya
da $!B$'nin geçerli olduğunu kanıtlamak da bunlardan hangisi olduğunu
kanıtlayabilmenizi gerektirir.

Biçimsel olarak sezgici birinci dereceden mantık, birinci dereceden mantık için
!!a{derivation} sistemini belirli bir biçimde kısıtlayınca elde edilen şeydir.
Benzer biçimde ikinci dereceden ya da yüksek dereceden mantığın sezgici
sürümleri de vardır. Matematiksel bakımdan bunlar yalnızca biçimsel tümdengelim
sistemleridir; ama daha önce belirtildiği gibi belirli bir matematiksel akıl
yürütme türünü modellemeleri amaçlanır. Bu, matematik hakkındaki belirli bir
felsefi görüşün (Brouwer'ın sezgiciliği gibi) gerekçelendirdiği akıl yürütme
türü olarak görülebilir; daha ``somut'' ve doyurucu bir matematiksel akıl
yürütme türü (Bishop'ın yapıcılığı doğrultusunda) sayılabilir; biçimsel
betimlemenin biçimsel olmayan güdüyü yakalayıp yakalamadığı tartışılabilir.
Fakat hangi felsefi görüşü benimsersek benimseyelim, sezgici mantığı biçimsel
olarak sunulmuş bir mantık diye inceleyebiliriz; her ne sebeple olursa olsun,
birçok matematiksel mantıkçı bunu yapmayı ilginç bulur.

Sezgici bağlaçların genellikle BHK yorumu (Brouwer, Heyting ve Kolmogorov'un
adlarından) diye bilinen, biçimsel olmayan yapıcı bir yorumu vardır. Yorum
şöyledir: $!A \land !B$'nin bir kanıtı, $!A$'nın bir kanıtıyla eşleştirilmiş
$!B$'nin bir kanıtından oluşur; $!A \lor !B$'nin bir kanıtı, hangisinin söz
konusu olduğuna ilişkin açık bilgiyle birlikte ya $!A$'nın ya da $!B$'nin bir
kanıtından oluşur; $!A \lif !B$'nin bir kanıtı, $!A$'nın bir kanıtını
$!B$'nin bir kanıtına dönüştüren bir yordamdan oluşur; $\lforall[x][!A(x)]$'in
bir kanıtı, bir $!A(x)$ kanıtını $x$'in her değeri için döndüren bir yordamdan
oluşur; $\lexists[x][!A(x)]$'in bir kanıtıysa, $x$'in bir değeri ile bu değerin
$!A$ koşulunu sağladığının bir kanıtından oluşur. Yorum, ``Curry--Howard
eşbiçimliliği'' ya da ``tipler-olarak-!!{formula}s paradigması'' adıyla bilinen
hesaplamalı terimlerle de betimlenebilir: !!a{formula}{}ün belirli bir veri
tipini belirttiğini ve kanıtların, karşılık gelen !!{formula}{}ün doğru olduğunu
görmemizi sağlayan bu veri tiplerinin hesaplamalı nesneleri olduğunu düşünün.

Sezgici mantık çoğu kez ``üçüncü hâlin olanaksızlığı yasası çıkarılmış'' klasik
mantık olarak düşünülür. Aşağıdaki teorem bunu daha kesin kılar.
\begin{thm}
Sezgici bakımdan aşağıdaki aksiyom şemaları eşdeğerdir:
\begin{enumerate}
\item $(\lnot !A \lif \lfalse) \lif !A$.
\item $!A \lor \lnot !A$
\item $\lnot \lnot !A \lif !A$
\end{enumerate}
\end{thm}
Şemalardan birinin örneklerini diğer ikisinden herhangi birinden elde etmek,
sezgici mantıkta iyi bir alıştırmadır.

Brouwer'ın sezgiciliğinin biçimselleştirmeleri olarak ortaya konan ilk sezgici
önermeler mantığı tümdengelim sistemleri, birbirinden bağımsız olarak
Kolmogorov, Glivenko ve Heyting'e aittir. Sezgici birinci dereceden mantığın
(ve sezgici matematiğin bazı bölümlerinin) ilk biçimselleştirmesi Heyting'e
aittir. Klasik bakımdan geçerli birçok şema sezgici bakımdan geçerli olmasa da
birçoğu geçerlidir.

\emph{Çifte değilleme çevirisi}, klasik mantıkla sezgici mantık arasındaki
önemli bir ilişkiyi betimler. Tümevarımlı olarak şöyle tanımlanır ($!A^N$'yi,
klasik $!A$ !!{formula}{}ünün ``sezgici'' çevirisi olarak düşünün):
\begin{align*}
!A^N & \ident \lnot\lnot !A \quad \text{atomik $!A$ !!{formula}{}ü için} \\
(!A \land !B)^N & \ident (!A^N \land !B^N) \\
(!A \lor !B)^N & \ident  \lnot\lnot (!A^N \lor !B^N) \\
(!A \lif !B)^N & \ident (!A^N \lif !B^N) \\
(\lforall[x][!A])^N & \ident \lforall[x][!A^N] \\
(\lexists[x][!A])^N & \ident \lnot\lnot\lexists[x][!A^N]
\end{align*}
Kolmogorov ile Glivenko'nun bu çevirinin önermeler mantığı için sürümleri
vardı; yüklem mantığı içinse çeviri birbirinden bağımsız olarak G\"odel ile
Gentzen'e aittir. Şunlar geçerlidir:

\begin{thm}
\begin{enumerate}
\item $!A \liff !A^N$ klasik olarak kanıtlanabilir.
\item $!A$ klasik olarak kanıtlanabiliyorsa, $!A^N$ sezgici olarak
  kanıtlanabilir.
\end{enumerate}
\end{thm}

Şimdi şu diyaloğu gözümüzde canlandırabiliriz. Klasik matematikçi: ``$!A$'yı
kanıtladım!'' Sezgici matematikçi: ``Kanıtın geçerli değil. Aslında kanıtladığın
$!A^N$.'' Klasik matematikçi: ``Benim için fark etmez!'' Klasik matematikçiye
göre sezgici matematikçi, ikisi eşdeğer olduğundan kılı kırk yarmaktadır. Ama
sezgici matematikçi arada bir fark bulunduğunu ısrarla savunur.

Yukarıdaki çevirinin yalnızca saf mantıkla ilgili olduğuna dikkat edin; klasik
ve sezgici matematik için uygun \emph{mantıksal olmayan} aksiyomların neler
olduğu ya da bunlar arasında nasıl bir ilişki bulunduğu sorusunu ele almaz.
Ancak yukarıdaki teoremin şu küçük genişletmesi bazı yararlı bilgiler sağlar:

\begin{thm}
$\Gamma$, $!A$'yı klasik olarak kanıtlıyorsa, $\Gamma^N$, $!A^N$'yi sezgici
olarak kanıtlar.
\end{thm}

Başka bir deyişle, $!A$ bazı varsayımlardan klasik olarak kanıtlanabiliyorsa,
$!A^N$ bu varsayımların çifte değilleme çevirilerinden kanıtlanabilir.

Bir tümcenin ya da önermesel !!{formula}{}ün sezgici bakımdan geçerli olduğunu
göstermek için yalnızca bir kanıt vermeniz yeterlidir. Peki geçerli
olmadığını nasıl gösterebilirsiniz? Bunun için sağlam ve tercihen tam
bir semantiğe gereksinim vardır. Kripke'ye ait bir semantik tam bu işe uyar.

Klasik mantıkta oynadığımız oyunun aynısını oynayabiliriz: semantiği tanımlar,
sağlamlık ile tamlığı kanıtlarız. Yine de şu ayrımı belirtmek yararlı olur.
Klasik mantıkta semantik, bir bakıma bağlaçların anlamında örtük bulunan
``apaçık'' semantikti. Kripke semantiği için bazı sezgisel gerekçeler
sunulabilse de bu semantik aynı kaçınılmazlık duygusunu vermez. Ayrıca klasik
!!{structure} kavramı doğal bir matematiksel kavramdır; dolayısıyla
!!a{structure} kavramını klasik birinci dereceden mantığı incelemekte bir araç,
ya da klasik birinci dereceden mantığı matematiksel !!{structure}{}sı incelemekte
bir araç sayabiliriz. Buna karşılık Kripke !!{structure}{}sı yalnızca mantıksal
bir kurgu olarak görülebilir; bağımsız bir matematiksel ilgi taşıyor
görünmezler.

Bir önermeler dili için bir Kripke !!{structure}{}sı
$\mModel M = \tuple{W, R, V}$, bir $W$ kümesinden, bir $R$ kısmi
sıralamasından (bu sıralamanın $W$ üzerinde en küçük bir !!{element}{}ı vardır) ve
önermesel değişkenleri $W$'nin !!{element}s{}ına atayan ``monoton'' bir atamadan
oluşur. Buradaki sezgi
şudur: $W$'nin !!{element}s{}ı ``dünyaları'' ya da ``bilgi durumlarını'' temsil
eder; $v \geq u$ elemanı $u$'nun ``olası bir gelecek durumunu'' temsil eder;
$u$'ya atanan önermesel değişkenler ise $u$ durumunda doğru olduğu bilinen
önermelerdir. Zorlama bağıntısı $\mSat{M}{!A}[w]$, daha sonra bu ilişkiyi
dildeki keyfî !!{formula}s için genişletir; $\mSat{M}{!A}[w]$ ifadesini
``$!A$, $w$ durumunda doğrudur'' diye okuyun. İlişki tümevarımlı olarak şöyle
tanımlanır:
\begin{enumerate}
\item $\mSat{M}{\Obj p_i}[w]$ ancak ve ancak $\Obj p_i$, $w$'ye atanan
  önermesel değişkenlerden biriyse geçerlidir.
\item $\mSat/{M}{\lfalse}[w]$.
\item $\mSat{M}{(!A \land !B)}[w]$ ancak ve ancak $\mSat{M}{!A}[w]$ ve
  $\mSat{M}{!B}[w]$ ise geçerlidir.
\item $\mSat{M}{(!A \lor !B)}[w]$ ancak ve ancak $\mSat{M}{!A}[w]$ ya da
  $\mSat{M}{!B}[w]$ ise geçerlidir.
\item $\mSat{M}{(!A \lif !B)}[w]$ ancak ve ancak $w' \geq w$ ve
  $\mSat{M}{!A}[w']$ olduğunda her zaman $\mSat{M}{!B}[w']$ ise geçerlidir.
\end{enumerate}
Bir karşı örnek sağlayan Kripke !!{structure}{}sı kurarak $\lnot (p \land q) \lif
(\lnot p \lor \lnot q)$ ifadesinin sezgici bakımdan geçerli olmadığını
göstermeye çalışmak iyi bir alıştırmadır.

